%%%%%
%
\section{Moyenne-Echantillonnage d'un param�tre \texttt{float}}

\itemstyle{Op�rateur~:} \operfmt{MoyEch}
\index{MoyEch}

%%%%%OBJET%%%%
\itemstyle{Objet~:}
Moyenne glissante et �chantillonnage d'un
param�tre � une fr�quence donn�e, diviseur
de la fr�quence du param�tre d'entr�e.

%%%%%ARGUMENTS%%%%
\bigskip
%%%%%ENTREES%%%%
\begin{tabarg}
\multicolumn{3}{l}{\titretab {Entr\'ees}}\\
\hline
$X$ & Vecteur, \texttt{float} & Param\`etre d'entr\'ee\\
$R_f$ & Coeff., \texttt{double} & Rapport des fr\'equences ($ F_{sortie}\over
F_{entree} $), dans l'intervalle $]0, 1]$ \\
$Decal$ & Coeff., \texttt{int} & D\'ecalage (\textit{nombre de points}),
dans l'intervalle $[-1/R_f, 1/R_f]$\\
\end{tabarg}

\medskip
%%%%%SORTIES%%%%
\begin{tabarg}
\multicolumn{3}{l}{\titretab {Sorties}}\\
\hline
$X_m$ & Vecteur, \texttt{float} & S\'erie des moyennes (m\^eme unit\'e que $X$)\\
\end{tabarg}


%%%%%FORMULE%%%%
\itemstyle{Formule}

%%MoyEch
\[
Xm[i] = R_f . \sum_{j = i/R_f + Decal}^{(i+1)/R_f -1 + Decal} X[j]
\qquad i=0,Nbre\_Echant(X). R_f
\]


%%%%COMMENTAIRE%%%%
\itemstyle{Commentaire~:}

$R_f$ doit �tre l'inverse d'un entier. Si ce n'est pas le cas, utiliser
\hyperref[OP-MoyEch_ff]{\operfmt{MoyEch\_ff}}.\newline
Le nombre d'�chantillons du vecteur de sortie est le produit du
nombre d'�chantillons du vecteur d'entr�e par $R_f$, arrondi
� l'entier sup�rieur. \newline
%Pour chaque �chantillon du r�sultat, s'il existe, dans la s�rie de
%donn�e en entr�e un seul �chantillon valeur manquantes 

%% Formulation a trouver pour les angles....
%%\emph{Attention} : Le calcul de moyenne convient pas aux variables
%%

%%%%%AUTEUR%%%%
\itemstyle{Auteur~:}
CNRM/GMEI/TRAMM

%%%%%REFERENCE%%%%
\itemstyle{R�f�rences~:}
